\documentclass[12pt, a4paper]{article}

%% --- Paquetes ---
\usepackage[utf8]{inputenc}
\usepackage[T1]{fontenc}
\usepackage[spanish]{babel}
\usepackage{geometry}
\usepackage{booktabs}
\usepackage{array}
\usepackage{setspace}
\usepackage{parskip}
\usepackage{titlesec}


\geometry{top=2.5cm, bottom=2.5cm, left=3cm, right=2.5cm}


\pagestyle{plain}

\titleformat{\section}{\large\bfseries}{}{0em}{}
\titleformat{\subsection}{\normalsize\bfseries}{}{0em}{}

\begin{document}


\begin{titlepage}
    \centering
    \vspace*{0.5cm}

    {\Large\bfseries
        Universidad Técnica Estatal de Quevedo\\[0.4cm]}

    {\normalsize Facultad de Ciencias de la Computación y Diseño Gráfico\\[0.1cm]
     Carrera de Ingeniería en Software\\[2cm]}

    \rule{\linewidth}{1.5pt}\\[0.4cm]
    {\LARGE\bfseries
        Sistema de Gestión de Solicitudes\\[0.2cm]
        para Soporte Técnico de Internet}\\[0.4cm]
    \rule{\linewidth}{1.5pt}\\[1.5cm]

    {\large\bfseries Propuesta de Proyecto para Aplicaciones Distribuidas\\[2cm]}

    %% --- Integrantes como lista ---
    \begin{flushleft}
        \textbf{Equipo de trabajo:}\\[0.4cm]
        \hspace{1cm}
        \begin{minipage}{0.85\textwidth}
            \begin{itemize}
                \setlength{\itemsep}{0.3cm}
                \item Pacheco Cárdenas Cristhian Daniel \quad --- \quad Arquitecto
                \item Carpio Mendoza Carlos José \quad --- \quad Líder de Desarrollo
                \item Álvarez Párraga Jeremy Alexis \quad --- \quad Responsable de Calidad
                \item Cando Moreno Robinson Rodrigo \quad --- \quad Responsable de Documentación
            \end{itemize}
        \end{minipage}
    \end{flushleft}

    \vspace{1.5cm}

    %% --- Datos del curso ---
    \begin{flushleft}
        \hspace{1cm}
        \begin{tabular}{ll}
            \textbf{Docente:} & Gleiston Cicerón Guerrero Ulloa \\[0.3cm]
            \textbf{Período Académico:} & Séptimo Semestre \\[0.3cm]
            \textbf{Fecha de Entrega:} & 27 de Mayo 2026 \\
        \end{tabular}
    \end{flushleft}

    \vfill
    {\small Quevedo -- Ecuador \quad | \quad 2026}
\end{titlepage}


\newpage
\onehalfspacing

\section{Resumen Ejecutivo}

\subsection{Problema que se resuelve}

Muchas organizaciones, especialmente de tamaño pequeño y mediano, gestionan
sus solicitudes de soporte técnico a través de canales informales como correos
electrónicos, llamadas telefónicas o mensajería instantánea. Esta forma de
trabajar genera desorganización, pérdida de información y tiempos de respuesta
prolongados. El presente proyecto propone un \textbf{Sistema de Gestión de
Solicitudes para Soporte Técnico de Internet}, orientado a formalizar y
estructurar el proceso de atención de incidencias relacionadas con la
conectividad a Internet dentro de entornos organizacionales.

\vspace{0.5cm}

\subsection{Usuarios del sistema}

El sistema está dirigido a organizaciones que ofrecen o dependen de servicios
de Internet y que necesitan un mecanismo ordenado para gestionar incidencias
técnicas. Lo utilizarán tanto el cliente que reporta problemas como el equipo
técnico encargado de resolverlos y el personal administrativo que supervisa y
coordina el proceso de atención.

\vspace{0.5cm}

\subsection{Un sistema distribuido es la solución apropiada}

Si esta aplicación se construyera como un sistema único y centralizado,
todo dependería de un solo punto: si ese punto falla, el servicio deja de
funcionar por completo.

Un enfoque distribuido resuelve esto dividiendo la aplicación en partes
separadas: la interfaz que ve el usuario, la lógica que procesa las
solicitudes y la base de datos donde se guarda la información. Cada parte
trabaja de forma independiente y se comunican entre sí. Esto trae beneficios
directos: si una parte falla, las demás pueden seguir funcionando; si la
cantidad de usuarios crece, se puede reforzar solo la parte que lo necesita;
y si se requiere hacer un cambio, se modifica únicamente la parte
correspondiente sin afectar el resto del sistema.

\newpage

\subsection{Alcance del período académico}

Durante el período académico, el trabajo consistirá en tomar la aplicación web del Sistema de Gestión de Solicitudes para Soporte Técnico de Internet, ya existente, e identificar de qué manera puede reorganizarse bajo un enfoque distribuido. A partir de ese análisis, se irán aplicando gradualmente los conceptos trabajados a lo largo del semestre sobre su estructura y diseño: cómo se comunican sus componentes, cómo se organiza y distribuye su información, y de qué forma puede estructurarse en capas independientes.

La idea es transformar la aplicación web original en un ecosistema basado en microservicios desacoplados y orientados a eventos, dividiendo la plataforma en componentes individuales de autenticación, gestión de reportes, notificaciones e inteligencia artificial que interactuarán de forma que no esperarse el uno al otro para seguir trabajando.

\end{document}
